\subsection{Actual proper valuative centers and bounded model stalks}
\label{subsec:proper-native-centers}

Let $R$ be any commutative ring and fix a real multiplicative ultrametric
seminorm $v$ with the explicit boundedness hypothesis $v(r)\leq1$ for every
$r\in R$. Let $f:X\to\operatorname{Spec}R$ be an actual proper scheme
morphism and let $y$ be a native valued point lying over the exact prescribed
base point. Write $x$ for its underlying scheme point and $w$ for its
valuation on the actual residue field $\kappa(x)$.
The native kernel support $x$ and the model center constructed below have
different definitions. No anti-continuity, analytic structure sheaf or
Noetherian hypothesis is asserted or used.

\begin{lemma}[Actual local residue-section pullback]
\label{scheme:actual-residue-section-pullback}
\lean{EnvelopeProperty.Berkovich.SchemeValuedPoint.fromSpecResidueField_appTop,
EnvelopeProperty.Berkovich.SchemeValuedPoint.fromSpecStalk_preimage_eq_top,
EnvelopeProperty.Berkovich.SchemeValuedPoint.fromSpecStalk_appLE_top,
EnvelopeProperty.Berkovich.SchemeValuedPoint.fromSpecResidueField_preimage_eq_top,
EnvelopeProperty.Berkovich.SchemeValuedPoint.fromSpecResidueField_appLE_top}
\leanok
\uses{scheme:valued-points}
The actual morphism $\operatorname{Spec}\kappa(x)\to X$ pulls every
section near $x$ back to its original residue value. The analogous map
from the spectrum of the local stalk pulls back to its genuine germ.
\end{lemma}
\begin{proof}
Every point of the spectrum of a local ring specializes to its closed point.
Consequently an open neighborhood of that point pulls back to the whole
local spectrum. The actual stalk-to-scheme section formula therefore reduces
to the germ map followed by the inverse global-section isomorphism of Spec.
The genuine residue map from the stalk, composition of actual section maps,
and Spec global-section naturality identify residue-field pullback with the
original germ followed by the residue map, exactly local residue evaluation.
This calculation preserves all sections, not only coefficient sections.
\end{proof}

\begin{definition}[Genuine bounded valuative square]
\label{scheme:bounded-native-valuative-square}
\lean{EnvelopeProperty.Berkovich.RelativeSchemeValuedPoint.residueBaseMap,
EnvelopeProperty.Berkovich.RelativeSchemeValuedPoint.residueBaseMap_val,
EnvelopeProperty.Berkovich.RelativeSchemeValuedPoint.CenterRing,
EnvelopeProperty.Berkovich.RelativeSchemeValuedPoint.centerBaseMap,
EnvelopeProperty.Berkovich.RelativeSchemeValuedPoint.centerBaseMap_coe,
EnvelopeProperty.Berkovich.RelativeSchemeValuedPoint.center_commSq,
EnvelopeProperty.Berkovich.RelativeSchemeValuedPoint.centerSquare}
\leanok
\uses{scheme:actual-residue-section-pullback,scheme:relative-valued-points,
scheme:relative-section-coefficients}
Let $V_y=\{a\in\kappa(x):w(a)\leq1\}$ be the actual valuation subring.
The actual structure morphism gives a coefficient residue map
$R\to\kappa(x)$ with valuation precisely $v$. The boundedness hypothesis
therefore makes it an actual map $R\to V_y$ and gives a genuine commutative
square
\[
\begin{array}{ccc}
\operatorname{Spec}\kappa(x)&\longrightarrow&X\\
\downarrow&&\downarrow f\\
\operatorname{Spec}V_y&\longrightarrow&\operatorname{Spec}R.
\end{array}
\]
The upper map is the genuine residue-field morphism, and
$\kappa(x)$ is the actual fraction field of the valuation ring $V_y$.
\end{definition}
\begin{proof}
Compose the genuine coefficient action on global sections with actual
residue evaluation. Prescribed-base section evaluation proves that each
coefficient has value $v(r)$. Thus $v(r)\leq1$ proves actual membership in
$V_y$, and restricting the ring-map codomain defines the lower morphism.
For the whole square, compare genuine maps to Spec by their global section
maps. The preceding residue pullback formula and the actual Spec global
section isomorphisms show equality on every coefficient. Mathlib's actual
valuation-subring structure supplies valuation-ring and fraction-field
instances, so this is a genuine valuative square without an assumed lift.
\end{proof}

\begin{theorem}[Actual unique proper center]
\label{scheme:proper-native-center}
\lean{EnvelopeProperty.Berkovich.RelativeSchemeValuedPoint.centerLiftStruct,
EnvelopeProperty.Berkovich.RelativeSchemeValuedPoint.centerLift,
EnvelopeProperty.Berkovich.RelativeSchemeValuedPoint.centerLift_generic,
EnvelopeProperty.Berkovich.RelativeSchemeValuedPoint.centerLift_base,
EnvelopeProperty.Berkovich.RelativeSchemeValuedPoint.centerLift_unique,
EnvelopeProperty.Berkovich.RelativeSchemeValuedPoint.center,
EnvelopeProperty.Berkovich.RelativeSchemeValuedPoint.support_specializes_center,
EnvelopeProperty.Berkovich.RelativeSchemeValuedPoint.support_mem_of_center_mem}
\leanok
\uses{scheme:bounded-native-valuative-square}
Properness constructs an actual morphism $\operatorname{Spec}V_y\to X$
whose two triangles commute. It is unique among whole scheme morphisms
with those two triangles. Define $c(y)$ as its image of the closed point
of $\operatorname{Spec}V_y$. Then $x$ specializes to $c(y)$, so every
open neighborhood of the model center contains the original native support.
\end{theorem}
\begin{proof}
Apply Mathlib's proper valuative criterion to the genuine square above.
It supplies the whole lift and both commuting triangles. The uniqueness
part follows from separatedness of the proper structure morphism.
The generic point of the valuation ring is the image of the unique point
of $\operatorname{Spec}\kappa(x)$, whose image in $X$ is $x$ by the first
triangle. That generic point specializes to the actual closed point of
$V_y$. Continuity of the constructed scheme morphism proves
$x\rightsquigarrow c(y)$. Open sets are stable under generization, giving
the neighborhood assertion. No kernel-support identification with the
center is assumed.
\end{proof}

\begin{theorem}[Genuine bounded model-section and stalk valuations]
\label{scheme:proper-center-local-bounds}
\lean{EnvelopeProperty.Berkovich.RelativeSchemeValuedPoint.centerSectionMap,
EnvelopeProperty.Berkovich.RelativeSchemeValuedPoint.centerSectionMap_coe,
EnvelopeProperty.Berkovich.RelativeSchemeValuedPoint.eval_le_one_of_center_mem,
EnvelopeProperty.Berkovich.RelativeSchemeValuedPoint.isBounded_sectionSpectrum_of_center_mem,
EnvelopeProperty.Berkovich.RelativeSchemeValuedPoint.centerStalkMap,
EnvelopeProperty.Berkovich.RelativeSchemeValuedPoint.centerStalkVal,
EnvelopeProperty.Berkovich.RelativeSchemeValuedPoint.centerStalkVal_le_one,
EnvelopeProperty.Berkovich.RelativeSchemeValuedPoint.centerStalkVal_lt_one_iff,
EnvelopeProperty.Berkovich.RelativeSchemeValuedPoint.centerStalkMap_germ,
EnvelopeProperty.Berkovich.RelativeSchemeValuedPoint.centerStalkVal_germ}
\leanok
\uses{scheme:proper-native-center,scheme:actual-residue-section-pullback,
scheme:relative-section-coefficients,def:general-ideal-norm}
For every actual scheme open $U$ containing $c(y)$, lift pullback gives an
actual ring map $\Gamma(X,U)\to V_y$. Its inclusion into $\kappa(x)$
is exactly original residue evaluation at $x$. Thus every section on $U$
has native value at most one, and its actual section-spectrum point is bounded.
The genuine local scheme map also gives $\mathcal O_{X,c(y)}\to V_y$.
Pulling back $w$ gives a bounded center-stalk valuation, agrees with every
original section value on its germ, and has value strictly smaller than one
precisely on the genuine maximal ideal of $\mathcal O_{X,c(y)}$.
\end{theorem}
\begin{proof}
The inverse image of $U$ in $\operatorname{Spec}V_y$ contains its closed
point, hence is the whole local spectrum. The actual section pullback and
global Spec isomorphism therefore land in $V_y$. The generic triangle,
composition of actual section maps, and the first lemma identify its inclusion
with original residue evaluation on every section. Membership in $V_y$
proves the section bound and actual spectrum boundedness.

The true local map at the closed point is Mathlib's stalk-closed-point map
of the constructed lift. It lands in $V_y$, so the induced stalk valuation
is bounded. The genuine germ/stalk pullback identity identifies its value
on each section germ with the preceding section value. The valuation ring
maximal ideal consists exactly of elements of value smaller than one; the
local-homomorphism property of the scheme stalk map pulls that maximal ideal
back to the actual center-stalk maximal ideal, proving the last equivalence.
Merely deriving an $R$-algebra action on an arbitrary section ring would not
provide these bounds; the actual proper valuative lift does.
\end{proof}
